\documentclass{tmsce}

% ---------------- Your packages ----------------
\usepackage{lipsum}
\usepackage{amsmath,amsfonts,amsthm}
\usepackage{algorithmic}
\usepackage{algorithm}
\usepackage{array}
\usepackage[caption=false,font=normalsize,labelfont=sf,textfont=sf]{subfig}
\usepackage{textcomp}
\usepackage{stfloats}
\usepackage{url}
\usepackage{verbatim}
\usepackage{graphicx}
\usepackage{cite}
\usepackage{mwe} % provides example-image (remove later when you use your real images)

\hyphenation{op-tical net-works semi-conduc-tor IEEE-Xplore}

% ---------------- Front-matter (keep blank for preprint / student paper) ----------------
\doi{}
\copyrightline{}
\permissions{}

\vol{}
\issue{}
\yearofpub{}
\pagerange{}

% ---------------- Title/Author ----------------
\title{Age- and Arm-Stratified Co-expression Network Rewiring in the Murine Kidney During and After Spaceflight}

\authors{Ibrahim M.~Shahid}
\affiliation{Middle College, The University of North Carolina at Greensboro, USA\\
\textit{Corresponding author:} imshahid@uncg.edu}

\keywords{spaceflight; kidney; transcriptomics; co-expression networks; network rewiring; aging}

% ---------------- Theorem environments (sample) ----------------
\newtheorem{theorem}{Theorem}
\newtheorem{lemma}{Lemma}

\begin{document}
% Remove journal-style footer/header for a preprint-style PDF.
\renewcommand{\journalname}{}
\makeatletter
\renewcommand{\tmsce@topstrip}{}
\makeatother
\fancyfoot[C]{}
\renewcommand{\footrulewidth}{0pt}
\maketitle

\begin{abstract}
Astronauts exhibit elevated nephrolithiasis risk following spaceflight, and prior work links microgravity to distal nephron remodeling and renal transporter dysregulation. Whether spaceflight alters the regulatory coordination among genes — shifting co-expression architecture without necessarily changing mean expression levels — and whether such rewiring differs by age or by spaceflight habitat, remains uncharacterized. We addressed this using OSD-771 (RRRM-2), bulk kidney RNA-seq from 80 female C57BL/6NTac mice under a full $2\times4\times2$ factorial design crossing Age (young, 16 weeks; old, 34 weeks), Environment Group (FLT, GC, VIV, BSL), and Arm (ISS Terminal sacrifice, ISS-T; Live Animal Return, LAR). Nephron segment proportions were estimated via MuSiC deconvolution against a murine single-cell kidney reference and used to remove compositional confounding while preserving biological contrasts. Sample-specific co-expression networks were inferred via LIONESS on a shared sparse skeleton constructed with cell-standardized Ledoit-Wolf partial correlations, and edge-wise regression over the full factorial design yielded model-predicted contrast-specific networks. Node2vec embeddings of these networks, aligned by orthogonal Procrustes across multi-seed runs, produced per-gene rewiring scores (cosine distance) for four primary flight contrasts. Focused permutation testing identified 11 significantly rewired genes in young ISS-T animals (FDR $< 0.05$), including the renal amino acid transporter Slc6a20a, the circadian regulator Cry2, and the vitamin D-binding protein Gc, with two additional hits in old ISS-T animals (Prlr, Nr4a1). Pathway enrichment revealed coordinated ECM and complement remodeling in young ISS-T mice, an EMT and steroid biosynthesis signature in old ISS-T mice, and retinoic acid metabolism as the most statistically robust signal across both arms — with strongest enrichment in the LAR recovery condition (OR = 19.9, $q = 1.6\times 10^{-5}$), suggesting active renal repair signaling during re-adaptation to terrestrial gravity. The ISS-T arm consistently showed greater and more significant rewiring than LAR across all contrasts, consistent with partial network recovery following return to Earth. Pre-registered NCC-WNK axis genes did not emerge among top individual rewiring hits, though pathway-level retinol and circadian signals offer candidate upstream mechanisms for known DCT dysregulation. This framework identifies retinoic acid signaling, ECM remodeling, and circadian network disruption as age-stratified transcriptomic responses to spaceflight invisible to standard differential expression analysis, and establishes a statistically defensible methodology for network rewiring inference in small-n factorial space biology datasets.
\end{abstract}

\printkeywords
% \printdates

% optional separator line (single-column)
\tmsceendfrontmatter

% =========================================================
\section{Introduction}
\lipsum[1]

We will refer to a classic text using a sample citation \cite{knuth84}.
We also demonstrate inline math, e.g., $e^{i\pi}+1=0$, and a displayed equation:
\begin{equation}\label{eq:sample}
\sum_{k=1}^{n} k = \frac{n(n+1)}{2}.
\end{equation}

\subsection{A sample theorem}
\begin{theorem}[Toy result]
For every integer $n\ge 1$, the sum $\sum_{k=1}^{n} k$ is equal to $\frac{n(n+1)}{2}$.
\end{theorem}
\begin{proof}
The proof can be done by induction. The base case $n=1$ holds. Assume true for $n$.
Then $\sum_{k=1}^{n+1} k = \left(\sum_{k=1}^{n} k\right) + (n+1) = \frac{n(n+1)}{2} + (n+1)=\frac{(n+1)(n+2)}{2}$.
\end{proof}

% =========================================================
\section{Sample Figure(s)}
Figure~\ref{fig:single} shows a sample figure. Figure~\ref{fig:subfig} shows two subfigures.

\begin{figure}[!t]
\centering
\includegraphics[width=0.75\linewidth]{example-image}
\caption{A sample figure (replace \texttt{example-image} with your own file).}
\label{fig:single}
\end{figure}

\begin{figure}[!t]
\centering
\subfloat[First subfigure]{%
  \includegraphics[width=0.45\linewidth]{example-image-a}
  \label{fig:subA}
}\hfill
\subfloat[Second subfigure]{%
  \includegraphics[width=0.45\linewidth]{example-image-b}
  \label{fig:subB}
}
\caption{A sample subfigure layout.}
\label{fig:subfig}
\end{figure}

% =========================================================
\section{Sample Table}
Table~\ref{tab:example} shows a simple table with aligned columns.

\begin{table}[!t]
\centering
\caption{A sample table for demonstration.}
\label{tab:example}
\renewcommand{\arraystretch}{1.2}
\begin{tabular}{|c|c|c|}
\hline
\textbf{Method} & \textbf{Time (ms)} & \textbf{Accuracy (\%)} \\
\hline
Baseline & 12.5 & 91.2 \\
Proposed & 10.1 & 93.8 \\
Optimized & 8.7 & 94.4 \\
\hline
\end{tabular}
\end{table}

% =========================================================
\section{Sample Algorithm}
Algorithm~\ref{alg:sample} shows a basic algorithm environment.

\begin{algorithm}[!t]
\caption{Sample procedure (toy example).}
\label{alg:sample}
\begin{algorithmic}[1]
\REQUIRE Input array $A[1..n]$
\ENSURE \texttt{true} if $A$ is non-decreasing, else \texttt{false}
\FOR{$i \leftarrow 1$ to $n-1$}
  \IF{$A[i] > A[i+1]$}
    \RETURN \texttt{false}
  \ENDIF
\ENDFOR
\RETURN \texttt{true}
\end{algorithmic}
\end{algorithm}

% =========================================================
\section{Conclusion}
We included samples of an equation (Eq.~\ref{eq:sample}), figures (Fig.~\ref{fig:single}--\ref{fig:subfig}),
a table (Table~\ref{tab:example}), and an algorithm (Algorithm~\ref{alg:sample}) to serve as a template.

\bibliographystyle{plain}
\begin{thebibliography}{1}

\bibitem{knuth84}
Donald~E. Knuth.
\newblock \emph{The \TeX book}.
\newblock Addison-Wesley, 1984.

\end{thebibliography}

\end{document}
